% Erklaerungen zum Type-I-Precoding und Beamforming in OCUDU.
% Uebersetzen: pdflatex main.tex  (zweimal fuer das Inhaltsverzeichnis)
\documentclass[a4paper,11pt]{article}
\usepackage[utf8]{inputenc}
\usepackage[T1]{fontenc}
\usepackage[ngerman]{babel}
\usepackage{lmodern}
\usepackage{amsmath,amssymb}
\usepackage{booktabs}
\usepackage{geometry}
\geometry{margin=2.5cm}
\usepackage{xcolor}
\usepackage{listings}
\usepackage{hyperref}
\hypersetup{colorlinks=true,linkcolor=blue!50!black,urlcolor=blue!50!black}

% Quellcode-Darstellung.
\definecolor{codebg}{gray}{0.96}
\definecolor{codecomment}{rgb}{0.3,0.5,0.3}
\definecolor{codekw}{rgb}{0.1,0.1,0.6}
\lstset{
  basicstyle=\ttfamily\small,
  backgroundcolor=\color{codebg},
  commentstyle=\color{codecomment},
  keywordstyle=\color{codekw}\bfseries,
  breaklines=true,
  columns=fullflexible,
  keepspaces=true,
  frame=single,
  framerule=0pt,
  xleftmargin=4pt,
  xrightmargin=4pt,
  showstringspaces=false,
  tabsize=2,
}
\lstdefinelanguage{yaml}{
  keywords={true,false,null},
  sensitive=false,
  comment=[l]{\#},
  morestring=[b]",
}
\newcommand{\file}[1]{\texttt{#1}}
\newcommand{\code}[1]{\texttt{#1}}

\title{Type-I-Precoding und Beamforming in OCUDU\\\large Erklärungen mit Quellcode-Ausschnitten}
\author{Zusammengestellt aus einer Code-Analyse des Repositories}
\date{\today}

\begin{document}
\maketitle
\tableofcontents
\bigskip

\noindent Alle Pfade sind relativ zur Wurzel des Repositories. Zeilennummern beziehen sich auf den Stand vom
6.\,November 2026 (Commit \code{9f591a95fc}) und können sich mit späteren Änderungen verschieben.

\section{Wo die Type-I-Codebooks erzeugt werden}

Die Type-I-Codebooks sind in zwei Schichten organisiert. Die Matrix-Mathematik liegt in \file{lib/ran/precoding/},
die Erzeugung des kompletten Codebook-Repositories für eine Zelle im FAPI-Adaptor.

\subsection{Die Mathematik: eine Matrix aus einem PMI}

\begin{itemize}
  \item \file{lib/ran/precoding/type1\_sp/precoding\_codebooks.cpp} ist der Kern. \code{make\_type1\_sp\_mode1()}
    baut die Gewichtsmatrix nach TS\,38.214 Tabellen 5.2.2.2.1-5 bis -8 für 1 bis 8 Layer. Pro Layer wird
    \code{make\_layer\_type1\_sp\_mode1()} aufgerufen: Kronecker-Produkt der horizontalen und vertikalen
    DFT-Vektoren, die zweite Polarisation erhält den Co-Phasing-Faktor $\varphi$.
    \code{get\_k1\_and\_k2\_2\_layers()} und \code{get\_k1\_and\_k2\_3\_4\_layers()} sind die Tabellen
    5.2.2.2.1-3 und -4 (Beam-Offset zwischen Layern über $i_{1,3}$).
  \item Dieselbe Datei enthält \code{calculate\_mimo\_matrix()} und \code{get\_beams\_from\_pmi()}, die zerlegte
    Variante für Topologien mit Beam-Raster (Abschnitt~\ref{sec:zwei-matrizen}).
  \item \file{lib/ran/precoding/two\_port/precoding\_codebooks.cpp}: Sonderfall 2 Ports, Tabelle 5.2.2.2.1-1,
    hart codiert.
  \item \file{lib/ran/precoding/precoding\_codebooks.cpp}: Dispatcher. \code{make\_precoding()} und
    \code{get\_mimo\_matrix\_from\_pmi()} machen ein \code{std::visit} über die PMI-Variante
    (Abschnitt~\ref{sec:dispatcher}).
  \item \file{lib/ran/precoding/type1\_sp/precoding\_codebook\_helpers.cpp}: Wertebereiche und Bitbreiten der
    PMI-Felder $i_{1,1}, i_{1,2}, i_{1,3}, i_2$ je Panel und Rang, TS\,38.212 Tabelle 6.3.1.1.2-1.
  \item \file{lib/ran/precoding/precoding\_codebook\_configuration.cpp}: \code{get\_single\_panel\_info()}, die
    Tabelle $N_1, N_2, O_1, O_2$ der unterstützten Panels $2\times1$, $2\times2$, $4\times1$.
\end{itemize}

Die Typen stehen in \file{include/ocudu/ran/precoding/}: \file{precoding\_matrix\_indicator.h} mit
\code{pmi\_typeI\_single\_panel}, \file{precoding\_codebook\_configuration.h} und
\file{precoding\_codebook\_type1\_helpers.h}.

\begin{lstlisting}[language=C++,title={include/ocudu/ran/precoding/precoding\_matrix\_indicator.h}]
struct pmi_typeI_single_panel {
  /// Single-panel topology configuration.
  pmi_codebook_typeI_single_panel panel_config;
  /// PMI parameter i_{1,1}. Values: {0,...,N_1*O_1 - 1}, at most {0,...,63}.
  uint8_t i_1_1;
  /// PMI parameter i_{1,2}. Only available for N_2 > 1 or v > 2.
  std::optional<uint8_t> i_1_2;
  /// PMI parameter i_{1,3}. Only available for v in {2,3,4}.
  std::optional<uint8_t> i_1_3;
  /// PMI parameter i_2.
  uint8_t i_2;
};

using precoding_matrix_indicator =
    std::variant<std::monostate, pmi_two_antenna_port, pmi_typeI_single_panel, pmi_typeII>;
\end{lstlisting}

\subsection{Das komplette Codebook für die Zelle}

\file{lib/fapi\_adaptor/precoding\_codebook\_generator.cpp}, Funktion \code{generate\_precoding\_codebooks()}.
Beim Start wird für die Zelle einmal jede PMI-Kombination durchgezählt. In \code{generate\_pdsch\_sp\_type1()}
laufen vier verschachtelte Schleifen über $i_{1,1}, i_{1,2}, i_{1,3}, i_2$, pro Layer-Anzahl 1 bis
\code{nof\_ports}, und jede Matrix landet unter einem Precoding-Matrix-Index im Repository.

\begin{lstlisting}[language=C++,title={lib/fapi\_adaptor/precoding\_codebook\_generator.cpp, generate\_pdsch\_sp\_type1()}]
pmi_typeI_single_panel_param_ranges param_ranges = get_pmi_ranges_typeI_single_panel(panel, nof_layers);

for (uint8_t i_1_1 = 0; i_1_1 != param_ranges.i_1_1; ++i_1_1) {
  for (uint8_t i_1_2 = 0; i_1_2 != param_ranges.i_1_2; ++i_1_2) {
    for (uint8_t i_1_3 = 0; i_1_3 != param_ranges.i_1_3; ++i_1_3) {
      for (uint8_t i_2 = 0; i_2 != param_ranges.i_2; ++i_2) {
        pmi_typeI_single_panel pmi = {.panel_config = panel,
                                      .i_1_1        = i_1_1,
                                      .i_1_2        = (param_ranges.i_1_2 > 0) ? std::optional(i_1_2) : std::nullopt,
                                      .i_1_3        = (param_ranges.i_1_3 > 0) ? std::optional(i_1_3) : std::nullopt,
                                      .i_2          = i_2};

        unsigned pm_index = base_offset + get_pdsch_single_panel_type1_precoding_matrix_index(param_ranges, pmi);
        repo_builder.add(pm_index, make_composite(pmi, nof_layers, topology));
      }
    }
  }
}
\end{lstlisting}

Daneben werden Identität, SSB, PDCCH, omnidirektionales PDSCH, CSI-RS und PRS eingetragen.
\code{make\_composite()} entscheidet anhand der Antennentopologie, ob die volle Port-Matrix oder die
Beam-Zerlegung gespeichert wird:

\begin{lstlisting}[language=C++,title={lib/fapi\_adaptor/precoding\_codebook\_generator.cpp, make\_composite()}]
static precoding_beamforming_composite
make_composite(const precoding_matrix_indicator& pmi, unsigned nof_layers, antenna_topology topology)
{
  if (has_beam_grid(topology)) {
    return get_mimo_matrix_from_pmi(pmi, nof_layers);
  }

  precoding_weight_matrix precoding = make_precoding(pmi, nof_layers);

  return {precoding, get_default_beam_list(precoding.get_nof_ports())};
}
\end{lstlisting}

Aufgerufen wird die Generierung in \file{lib/du/du\_high/o\_du\_high\_factory.cpp:105} (MAC-Seite, Mapper) und
\file{lib/du/du\_low/o\_du\_low\_factory.cpp:64} (PHY-Seite, Repository). Welches Codebook gilt, hängt nur von
der Anzahl der Sendeantennen ab:

\begin{lstlisting}[language=C++,title={lib/du/du\_low/o\_du\_low\_factory.cpp}]
switch (fapi_cfg.sectors[i].p7_config.carrier_cfg.num_tx_ant) {
  case 1: codebook_config = pmi_codebook_one_port{}; break;
  case 2: codebook_config = pmi_codebook_two_port{}; break;
  case 4:
    codebook_config =
        pmi_codebook_typeI_single_panel{pmi_codebook_single_panel_config::two_one, pmi_codebook_typeI_mode::one};
    break;
  case 8:
    codebook_config =
        pmi_codebook_typeI_single_panel{pmi_codebook_single_panel_config::four_one, pmi_codebook_typeI_mode::one};
    break;
}
\end{lstlisting}

\subsection{Konfiguration und Nutzung im Scheduler}

\begin{itemize}
  \item \file{lib/scheduler/config/csi\_helper.cpp}, \code{make\_type1\_codebook\_config()}: baut die
    RRC-CodebookConfig für das UE (Subset-Restriction-Bitmaps, Panel $2\times1$ oder $4\times1$).
  \item \file{lib/scheduler/ue\_context/ue\_channel\_state\_manager.cpp}: Default-PMI-Konfigurationen und
    Speicherung des vom UE gemeldeten PMI.
\end{itemize}

Zum Einstieg empfiehlt sich die Reihenfolge \code{make\_type1\_sp\_mode1()} $\rightarrow$
\code{generate\_pdsch\_sp\_type1()} $\rightarrow$ \code{calculate\_mimo\_matrix()}.

\section{OCUDU mit PDSCH-Traffic starten und die Generierung beobachten}

Es gibt zwei Wege: die komplette gNB im Testmodus mit Dummy-RU (kein Core, kein Funk, aber echter PDSCH-Traffic
für ein Phantom-UE), oder die Unit-Tests der Codebook-Generierung im Debugger.

\subsection{Bauen}

\begin{lstlisting}[language=bash]
cmake --build cmake-build-debug --target gnb precoding_codebook_generator_test precoding_codebooks_test -j4
\end{lstlisting}

\subsection{Konfiguration: Testmodus mit 4 Ports}

Vorlage ist \file{configs/testmode.yml}. Für Type-I-Codebooks sind mindestens 4 DL-Antennen nötig, sonst wird
nur das 2-Port- oder 1-Port-Codebook erzeugt.

\begin{lstlisting}[language=yaml,title={configs/testmode\_4x4.yml (Beispiel)}]
cu_cp:
  amf:
    no_core: true

ru_dummy:
  dl_processing_delay: 1

cell_cfg:
  dl_arfcn: 632628
  band: 78
  channel_bandwidth_MHz: 20
  common_scs: 30
  plmn: "00101"
  tac: 7
  pci: 1
  nof_antennas_dl: 4
  nof_antennas_ul: 4
  # Optional: Topologie mit Beam-Raster. Ohne Angabe wird four_ports abgeleitet.
  # tx_ant_topology: single_panel_two_one

test_mode:
  test_ue:
    rnti: 0x44
    nof_ues: 1
    pdsch_active: true
    pusch_active: false
    cqi: 15
    ri: 2          # 1..4, bestimmt die Layer-Anzahl
    i_1_1: 1       # PMI-Felder, die das Phantom-UE "meldet"
    i_1_3: 0
    i_2: 1

log:
  filename: /tmp/gnb.log
  all_level: warning
  mac_level: debug
  phy_level: debug
\end{lstlisting}

Das Phantom-UE meldet seine CSI aus genau diesen Werten. Das Packen passiert in
\file{lib/du/du\_high/test\_mode/mac\_test\_mode\_helpers.cpp}:

\begin{lstlisting}[language=C++,title={lib/du/du\_high/test\_mode/mac\_test\_mode\_helpers.cpp, fill\_pmi\_bits()}]
const auto* single_panel = std::get_if<pmi_codebook_typeI_single_panel>(&pmi_codebook);
if (single_panel == nullptr) {
  return;
}

// The fields that the panel topology and the rank do not report have a zero bit-width.
const pmi_typeI_single_panel_param_sizes sizes =
    get_pmi_sizes_typeI_single_panel(get_single_panel_info(single_panel->n1_n2), rank);
packed.push_back(test_ue_cfg.i_1_1, sizes.i_1_1);
packed.push_back(test_ue_cfg.i_1_2.value_or(0), sizes.i_1_2);
packed.push_back(test_ue_cfg.i_1_3.value_or(0), sizes.i_1_3);
packed.push_back(test_ue_cfg.i_2, sizes.i_2);
\end{lstlisting}

Starten:
\begin{lstlisting}[language=bash]
./cmake-build-debug/apps/gnb/gnb -c configs/testmode_4x4.yml
\end{lstlisting}

Die Topologie-Namen liefert \code{./cmake-build-debug/apps/gnb/gnb cell\_cfg -h}. Namen mit Präfix
\code{single\_panel\_} haben ein Beam-Raster.

\subsection{Breakpoints für die Generierung}

Die Codebook-Erzeugung loggt nichts und läuft genau einmal beim Start, bevor Traffic fließt.
\begin{itemize}
  \item \file{lib/fapi\_adaptor/precoding\_codebook\_generator.cpp:184}, die Zeile
    \code{repo\_builder.add(pm\_index, make\_composite(...))} in der vierfachen Schleife.
  \item \file{lib/ran/precoding/type1\_sp/precoding\_codebooks.cpp}, Einstieg von \code{make\_type1\_sp\_mode1()}
    und der \code{switch (nof\_layers)}.
  \item \file{lib/du/du\_low/o\_du\_low\_factory.cpp:64} und \file{lib/du/du\_high/o\_du\_high\_factory.cpp:105},
    die beiden Aufrufer.
\end{itemize}

Isolierter geht es über die Unit-Tests:
\begin{lstlisting}[language=bash]
./cmake-build-debug/tests/unittests/fapi_adaptor/precoding_codebook_generator_test --gtest_filter='*TypeI*'
./cmake-build-debug/tests/unittests/ran/precoding_codebooks_test
\end{lstlisting}

\subsection{Die Nutzung im laufenden Traffic}

Mit \code{mac\_level: debug} schreibt der Scheduler pro PDSCH-Grant die gewählte PMI ins Log
(\file{lib/scheduler/logging/scheduler\_result\_logger.cpp:627}, Format \code{ri=<layers>} gefolgt von den
PMI-Feldern). Mit \code{phy\_level: debug} protokolliert der PDSCH-Prozessor die angewendete
Precoding-Konfiguration.

\begin{lstlisting}[language=bash]
grep -E "PDSCH.*ri=" /tmp/gnb.log | head
\end{lstlisting}

Den Weg von der PMI im Scheduler zum Index im Repository zeigt \file{lib/fapi\_adaptor/precoding\_codebook\_mapper.cpp}.
Zur Laufzeit wird nichts neu gerechnet: Der PHY-Adaptor holt in
\file{lib/fapi\_adaptor/phy/p7/pdu\_translators/precoding.h:19} mit \code{get\_precoding\_config()} das fertige
Composite aus dem Repository.

\section{3GPP-Codebook und Beam-Codebook: zwei verschiedene $W$}
\label{sec:zwei-matrizen}

Es gibt zwei verschiedene \glqq$W$\grqq{} im Code. Die 3GPP-Matrix ist \code{make\_type1\_sp\_mode1()}.
\code{generate\_beam\_weights\_codebook()} ist \emph{nicht} die Umsetzung des 3GPP-Codebooks.

\subsection{Das 3GPP-Codebook}

\file{lib/ran/precoding/type1\_sp/precoding\_codebooks.cpp}, \code{make\_type1\_sp\_mode1()}, erzeugt aus dem PMI
die Matrix $W$ mit Dimension CSI-RS-Ports $\times$ Layer, genau wie in TS\,38.214 Tabelle 5.2.2.2.1-5 bis -8.
Pro Layer ein Vektor $[v_{l,m};\,\varphi_n v_{l,m}]$, normiert mit $1/\sqrt{\upsilon P}$.

\subsection{Das Beam-Codebook}

\file{lib/ran/beamforming/beam\_weights\_codebook\_generator.cpp} steht in keiner 3GPP-Spezifikation. Es ist eine
implementierungsspezifische Beam-Tabelle: alle DFT-Richtungen, die das Antennen-Array bilden kann, plus die
Einzelport-Beams. Das Konzept stammt aus O-RAN, wo die Radio Unit eine feste Tabelle von Beam-IDs kennt.
Zeilen sind Beams, Spalten physische Antennen.

\begin{lstlisting}[language=C++,title={lib/ran/beamforming/beam\_weights\_codebook\_generator.cpp}]
beam_weights_codebook ocudu::generate_beam_weights_codebook(antenna_topology topology)
{
  // ... Topologie-Parameter: nof_panels, nof_elements_dim1/2, nof_polarizations, nof_beams_dim1/2 ...
  beam_weights_codebook beam_weights = beam_weights_codebook(nof_total_beams, nof_total_antennas);

  // Direkte Port-Beams: Einheitsmatrix in den ersten nof_total_antennas Zeilen.
  for (unsigned i_antenna = 0; i_antenna != nof_total_antennas; ++i_antenna) {
    beam_weights.set_coefficient(1.0, get_beam_id(topology, i_antenna), i_antenna);
  }

  // Only direct beam to port mapping if the panels are 1x1 elements.
  if (!has_beam_grid(topology)) {
    return beam_weights;
  }

  float amplitude = std::sqrt(1.0F / static_cast<float>(nof_elements_dim1 * nof_elements_dim2));

  for (unsigned i_panel = 0; i_panel != nof_panels; ++i_panel) {
    for (unsigned i_pol = 0; i_pol != nof_polarizations; ++i_pol) {
      for (unsigned i_beam_dim1 = 0; i_beam_dim1 != nof_beams_dim1; ++i_beam_dim1) {
        for (unsigned i_beam_dim2 = 0; i_beam_dim2 != nof_beams_dim2; ++i_beam_dim2) {
          beam_identifier beam_id = get_beam_id(topology, i_panel, i_pol, i_beam_dim1, i_beam_dim2);
          unsigned i_port = nof_elements_dim1 * nof_elements_dim2 * (nof_polarizations * i_panel + i_pol);

          for (unsigned j = 0; j != nof_elements_dim1; ++j) {
            cf_t v_l = std::polar(amplitude, TWOPI * j * i_beam_dim1 / nof_beams_dim1);
            for (unsigned k = 0; k != nof_elements_dim2; ++k) {
              cf_t u_m   = std::polar(1.0F, TWOPI * k * i_beam_dim2 / nof_beams_dim2);
              beam_weights.set_coefficient(v_l * u_m, beam_id, i_port++);
            }
          }
        }
      }
    }
  }
  return beam_weights;
}
\end{lstlisting}

Ausgabe für \code{single\_panel\_two\_one} (20 Beams $\times$ 4 Antennen, gekürzt):
\begin{lstlisting}
  beam  0: [+1.00+0.00j  +0.00+0.00j  +0.00+0.00j  +0.00+0.00j]
  beam  3: [+0.00+0.00j  +0.00+0.00j  +0.00+0.00j  +1.00+0.00j]
  beam  4: [+0.71+0.00j  +0.71+0.00j  +0.00+0.00j  +0.00+0.00j]   l=0, pol 0
  beam  5: [+0.71+0.00j  +0.50+0.50j  +0.00+0.00j  +0.00+0.00j]   l=1, pol 0
  beam 12: [+0.00+0.00j  +0.00+0.00j  +0.71+0.00j  +0.71+0.00j]   l=0, pol 1
  beam 13: [+0.00+0.00j  +0.00+0.00j  +0.71+0.00j  +0.50+0.50j]   l=1, pol 1
\end{lstlisting}

\subsection{Die Zerlegung}

Die Steuervektoren $v_{l,m}$ des Type-I-Codebooks sind DFT-Vektoren mit Phase $2\pi l i/(N_1 O_1)$. Das
Beam-Codebook enthält dieselben Vektoren. Deshalb lässt sich jede 3GPP-Matrix zerlegen:
\begin{equation}
  W_{\text{3GPP}}\;(\text{Antennen}\times\text{Layer})
  = W_{\text{Beam}}^{\mathsf T}\;(\text{Antennen}\times\text{Beams})\cdot
    W_{\text{MIMO}}\;(\text{Beams}\times\text{Layer}).
\end{equation}
$W_{\text{MIMO}}$ ist winzig: pro Layer zwei Einträge, einer pro Polarisation mit $\pm\varphi$. Es entsteht in
\code{calculate\_mimo\_matrix()}:

\begin{lstlisting}[language=C++,title={lib/ran/precoding/type1\_sp/precoding\_codebooks.cpp, calculate\_mimo\_matrix()}]
precoding_beam_list beams = get_beams_from_pmi(pmi, nof_layers);   // Beam-IDs aus i_1_1, i_1_2, i_1_3
unsigned nof_beams = beams.size() / 2;                               // ohne Polarisationsdimension

unsigned i_2 = pmi.i_2 % pmi_ranges.i_2;
cf_t     phi = std::polar(1.0F, (TWOPI / 4.0F) * static_cast<float>(i_2));

precoding_weight_matrix weights(nof_layers, beams.size());

// Normalize precoding by the number of layers and ports, apply the reciprocal of the beamforming scaling.
float scaling = std::sqrt(static_cast<float>(info.n1 * info.n2)) / std::sqrt(static_cast<float>(nof_layers * nof_ports));

for (unsigned i_layer = 0; i_layer != nof_layers; ++i_layer) {
  unsigned i_beam    = i_layer % nof_beams;
  // First half of the layers uses +phi, second half uses -phi, as per TS38.214 Section 5.2.2.2.1.
  cf_t     layer_pol = (i_layer < (nof_layers + 1) / 2) ? phi : -phi;

  weights.set_coefficient(scaling,             i_layer, 2 * i_beam);      // Polarisation 0
  weights.set_coefficient(scaling * layer_pol, i_layer, 2 * i_beam + 1);  // Polarisation 1
}
return {weights, beams};
\end{lstlisting}

Der Precoding-Code kennt das Beam-Codebook nicht, nur die Beam-\emph{Nummern} über \code{get\_beam\_id()}.
Beide Seiten sind ausschließlich über diese Nummerierungskonvention gekoppelt. Dass sie zusammenpassen, sichert der
Test \file{precoding\_matrix\_factorization\_test}.

\paragraph{Zahlenbeispiel} \code{single\_panel\_two\_one}, 2 Layer, $i_{1,1}=1$, $i_{1,3}=0$, $i_2=1$:
beide Layer auf Beam $l=1$, Beam-IDs 5 (Pol.~0) und 13 (Pol.~1), $\varphi=j$, Skalierung
$\sqrt{2}/\sqrt{2\cdot4}=0{,}5$.
\begin{equation*}
  W_{\text{MIMO}} = \begin{bmatrix} 0{,}5 & +0{,}5j \\ 0{,}5 & -0{,}5j \end{bmatrix}
  \quad(\text{Layer}\times\text{Beams }5,13)
\end{equation*}
Layer 0 auf den Antennen: $0{,}5\cdot\text{beam}_5 + 0{,}5j\cdot\text{beam}_{13}
= [\,0{,}354,\; 0{,}25+0{,}25j,\; 0{,}354j,\; -0{,}25+0{,}25j\,]$, identisch mit der Spalte aus
\code{make\_type1\_sp\_mode1()}.

\subsection{Wo die Multiplikation mit den IQ-Samples passiert}

\begin{enumerate}
  \item \textbf{Upper PHY, Precoding.} Der Resource-Grid-Mapper (\file{lib/phy/support/resource\_grid\_mapper\_impl.cpp:127})
    ruft pro PRG \code{precoder.apply\_precoding()}. Mit Beam-Raster ist das $W_{\text{MIMO}}$, ohne die volle
    3GPP-Matrix. Die Rechnung steht in \file{lib/phy/generic\_functions/precoding/channel\_precoder\_generic.cpp}:
\begin{lstlisting}[language=C++]
for (unsigned i_re = 0; i_re != nof_re; ++i_re) {
  // Set the port RE to the contribution of the first layer.
  cf_t sum = layer_re_view_list[0][i_re] * port_weights[0];
  for (unsigned i_layer = 1; i_layer != nof_layers; ++i_layer) {
    // Accumulate the contributions of all other layers.
    sum += layer_re_view_list[i_layer][i_re] * port_weights[i_layer];
  }
  port_re[i_re] = sum;
}
\end{lstlisting}
  \item \textbf{Lower PHY, Beamforming.} Beim OFDM-Modulieren holt
    \file{lib/phy/lower/processors/downlink/pdxch/pdxch\_baseband\_modulator.h:163} für jede physische Antenne die
    Spalte des Beam-Codebooks und der Modulator bildet die gewichtete Summe über die Beam-Ports des Grids.
    Nullgewichte werden übersprungen (\code{std::isnormal}-Prüfung in
    \file{lib/phy/lower/modulation/ofdm\_modulator\_impl.cpp:86}).
\end{enumerate}

\subsection{Wann \code{generate\_beam\_weights\_codebook()} aufgerufen wird}

Genau einmal pro Zelle beim Start, unabhängig vom Traffic, in der Initialisierungsliste von
\code{pdxch\_processor\_impl} (\file{lib/phy/lower/processors/downlink/pdxch/pdxch\_processor\_impl.h:52}) sowie für
Open Fronthaul in \file{lib/ofh/transmitter/ofh\_transmitter\_factories.cpp:68}.

\begin{lstlisting}[language=C++,title={lib/phy/lower/processors/downlink/pdxch/pdxch\_processor\_impl.h}]
pdxch_processor_impl(std::unique_ptr<ofdm_symbol_modulator> modulator,
                     task_executor&                         executor,
                     const configuration&                   config) :
  logger(ocudulog::fetch_basic_logger("PHY")),
  beamforming_codebook(generate_beam_weights_codebook(config.tx_ant_topology)),
  ...
\end{lstlisting}

Ohne \code{tx\_ant\_topology} in der Konfiguration wird die Topologie aus \code{nof\_antennas\_dl} abgeleitet:
\code{one\_port} bis \code{eight\_ports}. Keine davon hat ein Beam-Raster, das Codebook ist dann die Einheitsmatrix
und wirkungslos.

\begin{table}[h]\centering
\begin{tabular}{lll}
\toprule
Setup & Beamforming geschieht in & Beam-Codebook \\
\midrule
SDR, \code{four\_ports} & Upper PHY, volle 3GPP-Matrix & Einheitsmatrix, wirkungslos \\
SDR, \code{single\_panel\_two\_one} & Upper PHY $W_{\text{MIMO}}$, Lower PHY Tabelle & nötig \\
OFH-RU mit Beamforming & Upper PHY $W_{\text{MIMO}}$, RU Tabelle & nötig, wird an RU gesendet \\
\bottomrule
\end{tabular}
\end{table}

\subsection{Wofür das Beam-Codebook trotzdem existiert}

\begin{itemize}
  \item \textbf{Open Fronthaul mit Beamforming in der Radio Unit.} Die DU sendet pro Beam-Port und nennt im
    C-Plane die Beam-ID. \file{lib/ofh/support/beamforming\_weights\_repository.cpp} komprimiert jede Zeile der
    Tabelle einmal und schickt sie an die RU.
  \item \textbf{Beams ohne PMI.} SSB, PDCCH, SIB und CSI-RS haben kein UE, das einen PMI meldet. Die
    Zellkonfiguration zeigt mit \code{ref\_beams} direkt auf Zeilen der Tabelle (z.\,B. SSB-Sweeping).
  \item \textbf{Ein Grid-Format für beide Backends.} Der Upper PHY schreibt immer dasselbe Resource Grid mit
    Beam-Ports, egal ob dahinter die eigene Lower PHY oder eine OFH-RU hängt.
\end{itemize}

\subsection{Der UHD-SDR-Fall}

Mit einem UHD-SDR und Standardtopologie passiert das komplette Beamforming im Upper PHY mit der 3GPP-Matrix:
\begin{enumerate}
  \item Codebook beim Start: \code{make\_type1\_sp\_mode1()} für jeden PMI, FAPI-Repository speichert.
  \item Precoding pro PDSCH: \code{apply\_precoding\_port()}, pro Port eine Spalte von $W$, pro Resource Element
    eine Summe über die Layer. Das ist $W\cdot x$ auf Subträgerebene, vor der IFFT.
  \item OFDM-Modulation: pro Port IFFT und Cyclic Prefix. Beam-Codebook ist Einheitsmatrix.
  \item UHD: die Antennen-Zeitsignale gehen als TX-Kanäle an das Radio. Die relative Phase zwischen den Kanälen ist
    die Keule; die Hardware muss phasenkohärent senden (gemeinsamer LO, z.\,B. N310/X310).
\end{enumerate}
Konfiguration: \code{nof\_antennas\_dl: 4}, \code{tx\_ant\_topology} nicht setzen. Das Codebook ist dann
\code{two\_one}, 2 kreuzpolarisierte Elemente, 8 Richtungen über $i_{1,1}$. Für 16 Richtungen mit 4 Elementen pro
Polarisation sind 8 TX-Kanäle nötig (\code{four\_one}).

\subsection{Was OCUDU beim Beamforming nicht macht}

\begin{itemize}
  \item Keine Keule auf einen beliebigen Winkel in Grad, nur das feste DFT-Raster.
  \item Keine Nebenkeulenunterdrückung durch Amplituden-Tapering; alle Koeffizienten haben denselben Betrag.
  \item Kein Downlink-Beamforming aus SRS-Messungen; die SRS-Kanalmatrix dient nur der Uplink-TPMI-Auswahl.
  \item Keine Nullstellen auf andere UEs, kein Zero-Forcing für Multi-User-MIMO.
\end{itemize}
Die Richtung bestimmt das UE über den PMI-Report; OCUDU wendet sie an. Im Testmodus setzt man sie mit $i_{1,1}$ fest.

\section{Die Mathematik von \code{make\_type1\_sp\_mode1()}}

TS\,38.214 Abschnitt 5.2.2.2.1, Type I Single Panel, Codebook-Mode 1.

\subsection{Antennenmodell}

Ein Panel mit $N_1$ Elementen horizontal, $N_2$ vertikal, jedes kreuzpolarisiert. CSI-RS-Ports $P = 2 N_1 N_2$.
Die ersten $N_1 N_2$ Ports sind Polarisation 0, die zweiten Polarisation 1. $O_1$, $O_2$ sind die
Oversampling-Faktoren; die Zahl der Richtungen ist $N_1 O_1$ horizontal und $N_2 O_2$ vertikal.

\begin{table}[h]\centering
\begin{tabular}{lcccccc}
\toprule
Codebook & $N_1$ & $N_2$ & $O_1$ & $O_2$ & Ports & Richtungen \\
\midrule
\code{two\_one}  & 2 & 1 & 4 & 1 & 4 & 8 \\
\code{two\_two}  & 2 & 2 & 4 & 4 & 8 & $8\times8$ \\
\code{four\_one} & 4 & 1 & 4 & 1 & 8 & 16 \\
\bottomrule
\end{tabular}
\end{table}

\subsection{Der Steuervektor $v_{l,m}$}

\begin{align}
  u_m &= \begin{bmatrix} 1 & e^{j\frac{2\pi m}{O_2 N_2}} & \cdots & e^{j\frac{2\pi m (N_2-1)}{O_2 N_2}} \end{bmatrix} \\
  v_{l,m} &= \begin{bmatrix} u_m & u_m\, e^{j\frac{2\pi l}{O_1 N_1}} & \cdots & u_m\, e^{j\frac{2\pi l (N_1-1)}{O_1 N_1}} \end{bmatrix}^{\mathsf T}
\end{align}
Ein Kronecker-Produkt aus horizontalem und vertikalem DFT-Vektor, Länge $N_1 N_2$. $l$ wählt die horizontale
Richtung, $m$ die vertikale. Alle Elemente haben Betrag 1, es gibt kein Tapering. Bei Elementabstand
$d=\lambda/2$ liegt die Hauptkeule bei $\sin\theta = 2l/(O_1 N_1)$ für $l \le N_1 O_1/2$, negative Winkel für die
obere Hälfte von $l$.

Im Code wird der Vektor durch fortlaufende Multiplikation mit einem Phasor aufgebaut:

\begin{lstlisting}[language=C++,title={lib/ran/precoding/type1\_sp/precoding\_codebooks.cpp, make\_layer\_type1\_sp\_mode1()}]
static void make_layer_type1_sp_mode1(precoding_weight_matrix&              result,
                                      const pmi_codebook_single_panel_info& panel,
                                      unsigned                              i_layer,
                                      unsigned                              l,
                                      unsigned                              m,
                                      cf_t                                  polarization_offset)
{
  // Phase increment for each beam coefficient. This defines the direction of the beam in the horizontal plane.
  float phase_inc_h_rad = TWOPI * (static_cast<float>(l) / static_cast<float>(panel.o1 * panel.n1));
  // Phase increment for each beam coefficient. This defines the direction of the beam in the vertical plane.
  float phase_inc_v_rad = TWOPI * (static_cast<float>(m) / static_cast<float>(panel.o2 * panel.n2));
  // Get the number of physical ports, which is half the number of CSI-RS ports given that two polarizations are used.
  unsigned nof_ports_2 = result.get_nof_ports() / 2;

  cf_t phase_inc_h = std::polar(1.0F, phase_inc_h_rad);
  cf_t phase_inc_v = std::polar(1.0F, phase_inc_v_rad);
  cf_t beam_coef_h = std::polar(1.0F, 0.0F);
  cf_t beam_coef_v = std::polar(1.0F, 0.0F);

  unsigned i_port = 0;
  for (unsigned i_h = 0; i_h != panel.n1; ++i_h) {
    beam_coef_v = std::polar(1.0F, 0.0F);
    for (unsigned i_v = 0; i_v != panel.n2; ++i_v) {
      cf_t beam_coef = beam_coef_h * beam_coef_v;
      // First polarization.
      result.set_coefficient(beam_coef, i_layer, i_port);
      // Second polarization, same beam but with phase offset.
      result.set_coefficient(polarization_offset * beam_coef, i_layer, nof_ports_2 + i_port);
      i_port++;
      beam_coef_v *= phase_inc_v;
    }
    beam_coef_h *= phase_inc_h;
  }
}
\end{lstlisting}

\subsection{Die Co-Phasierung $\varphi_n$}

\begin{equation}
  \varphi_n = e^{j\pi n/2},\qquad n = i_2 \in \{0,1,2,3\}\quad\Rightarrow\quad \varphi \in \{1,\, j,\, -1,\, -j\}
\end{equation}
Die Phase der zweiten Polarisation relativ zur ersten. Im Code \code{phi = std::polar(1, (TWOPI/4)*i\_2)}.

\subsection{Ein Layer}
\begin{equation}
  W^{(1)} = \frac{1}{\sqrt{P}}\begin{bmatrix} v_{l,m} \\ \varphi_n v_{l,m} \end{bmatrix}\qquad (P\times1),
  \quad l = i_{1,1},\; m = i_{1,2}
\end{equation}
Tabelle 5.2.2.2.1-5.

\subsection{Zwei Layer}
\begin{equation}
  W^{(2)} = \frac{1}{\sqrt{2P}}\begin{bmatrix} v_{l,m} & v_{l',m'} \\ \varphi_n v_{l,m} & -\varphi_n v_{l',m'} \end{bmatrix}
  \qquad (P\times2),\quad l' = l + k_1,\; m' = m + k_2
\end{equation}
Die Offsets $(k_1,k_2)$ kommen aus $i_{1,3}$ über Tabelle 5.2.2.2.1-3, im Code \code{get\_k1\_and\_k2\_2\_layers()}.
Für \code{two\_one} ist $k_1\in\{0, O_1\}$: derselbe Beam oder der nächste Hauptkeulen-Nachbar. Das Minus in der
zweiten Spalte macht die Layer orthogonal, auch bei $k_1=0$: $[v;\varphi v]\perp[v;-\varphi v]$. $i_2$ ist hier
nur noch 0 oder 1.

\subsection{Drei und vier Layer} ($P<16$)
\begin{align}
  W^{(3)} &= \frac{1}{\sqrt{3P}}\begin{bmatrix} v_{l,m} & v_{l',m'} & v_{l,m} \\ \varphi_n v_{l,m} & \varphi_n v_{l',m'} & -\varphi_n v_{l,m} \end{bmatrix} \\
  W^{(4)} &= \frac{1}{\sqrt{4P}}\begin{bmatrix} v_{l,m} & v_{l',m'} & v_{l,m} & v_{l',m'} \\ \varphi_n v_{l,m} & \varphi_n v_{l',m'} & -\varphi_n v_{l,m} & -\varphi_n v_{l',m'} \end{bmatrix}
\end{align}
$(k_1,k_2)$ nach Tabelle 5.2.2.2.1-4, \code{get\_k1\_and\_k2\_3\_4\_layers()}; hier ist $k=0$ nicht erlaubt.

\begin{lstlisting}[language=C++,title={make\_type1\_sp\_mode1(), Auswahl der Spalten}]
switch (nof_layers) {
  case 1: {
    layer_beam_horizontal = {i_1_1};
    layer_beam_vertical   = {i_1_2};
    layer_pol             = {phi};
    break;
  }
  case 2: {
    std::tie(k1, k2) = get_k1_and_k2_2_layers(panel_info, i_1_3);
    layer_beam_horizontal = {i_1_1, i_1_1 + k1};
    layer_beam_vertical   = {i_1_2, i_1_2 + k2};
    layer_pol             = {phi, -phi};
    break;
  }
  case 3:
  case 4: {
    std::tie(k1, k2)      = get_k1_and_k2_3_4_layers(panel_info, i_1_3);
    layer_beam_horizontal = {i_1_1, i_1_1 + k1, i_1_1, i_1_1 + k1};
    layer_beam_vertical   = {i_1_2, i_1_2 + k2, i_1_2, i_1_2 + k2};
    layer_pol             = {phi, phi, -phi, -phi};
    break;
  }
  // ... Faelle 5 bis 8: Tabellen 5.2.2.2.1-9 bis -12 ...
}

// Build the coefficients for each port and layer.
for (unsigned i_layer = 0; i_layer != nof_layers; ++i_layer) {
  make_layer_type1_sp_mode1(result, panel_info, i_layer,
                            layer_beam_horizontal[i_layer],
                            is_2d_panel ? layer_beam_vertical[i_layer] : 0,
                            layer_pol[i_layer]);
}

// Apply scaling.
float scaling = 1.0F / std::sqrt(nof_layers * nof_ports);
result *= scaling;
\end{lstlisting}

Fünf bis acht Layer benutzen bis zu vier Beams mit festen Offsets $O_1$ und $2O_1$ und abwechselnd $\varphi$ und
$\pm1$ (Tabellen 5.2.2.2.1-9 bis -12).

\subsection{Normierung}
\begin{equation}
  \text{scaling} = \frac{1}{\sqrt{\upsilon P}}
\end{equation}
$\upsilon$ Layer, $P$ Ports. Jeder Eintrag hat Betrag $1/\sqrt{\upsilon P}$, die Sendeleistung über alle Ports und
Layer ist 1.

\subsection{Zahlenbeispiel}

\code{two\_one}, 2 Layer, $i_{1,1}=1$, $i_{1,3}=0$, $i_2=1$:
\begin{equation*}
  v_1 = \begin{bmatrix} 1 \\ e^{j2\pi/8} \end{bmatrix} = \begin{bmatrix} 1 \\ 0{,}707+0{,}707j \end{bmatrix},\quad
  \varphi = j,\quad k_1 = 0
\end{equation*}
\begin{equation*}
  W^{(2)} = \frac{1}{\sqrt 8}\begin{bmatrix}
    1 & 1 \\
    0{,}707+0{,}707j & 0{,}707+0{,}707j \\
    j & -j \\
    -0{,}707+0{,}707j & 0{,}707-0{,}707j
  \end{bmatrix}
\end{equation*}
Spalte 1 ist $[\,0{,}354,\; 0{,}25+0{,}25j,\; 0{,}354j,\; -0{,}25+0{,}25j\,]$, derselbe Vektor wie aus der
Zerlegung in Abschnitt~\ref{sec:zwei-matrizen}.

\subsection{Was das Precoding damit macht}
\begin{equation}
  y_p[k] = \sum_{\upsilon} W[p,\upsilon]\, x_\upsilon[k]
  \qquad\text{für jeden Subträger } k \text{ und jeden Port } p
\end{equation}
Das ist die Schleife in \code{apply\_precoding\_port()}: pro Port eine Zeile von $W$, pro Resource Element die
Summe über die Layer. Am Empfänger in Richtung $\theta$ ergibt sich der Array-Faktor
\begin{equation}
  AF(\theta) = \sum_{i=0}^{N_1-1} e^{j2\pi i\left(\frac{l}{O_1 N_1} - \frac{d\sin\theta}{\lambda}\right)},
\end{equation}
maximal, wenn die Klammer null ist. Das ist die Hauptkeule; Nullstellen und Nebenkeulen folgen aus der Summe über
$N_1$ Elemente.

\section{Die beiden Dateien \file{precoding\_codebooks.cpp}}
\label{sec:dispatcher}

Zwei Dateien gleichen Namens mit klar getrennten Rollen: oben die Fassade und die codebook-unabhängigen
Matrizen, unten die eigentliche 3GPP-Type-I-Mathematik.

\subsection{\file{lib/ran/precoding/precoding\_codebooks.cpp}, die Fassade}

Implementiert den öffentlichen Header \file{include/ocudu/ran/precoding/precoding\_codebooks.h}.
\begin{itemize}
  \item Einfache Matrizen ohne PMI: \code{make\_single\_port()}, \code{make\_one\_layer\_one\_port()},
    \code{make\_one\_layer\_all\_ports()}, \code{make\_identity()}. Sie braucht der FAPI-Adaptor für SSB, PDCCH,
    CSI-RS und das omnidirektionale PDSCH.
  \item \code{get\_single\_panel\_topology()}: Abbildung von $2\times1$, $2\times2$, $4\times1$ auf die
    Antennentopologie.
  \item Zwei Dispatcher: \code{make\_precoding()} und \code{get\_mimo\_matrix\_from\_pmi()}. Die PMI ist ein
    \code{std::variant}; die Visitor-Structs rufen die passende Implementierung.
\end{itemize}

\begin{lstlisting}[language=C++,title={lib/ran/precoding/precoding\_codebooks.cpp}]
/// Precoding weight matrix calculator from a given PMI.
struct precoding_calculator {
  unsigned nof_layers;

  precoding_weight_matrix operator()(const pmi_two_antenna_port& pmi) const
  {
    switch (nof_layers) {
      case 1: return make_one_layer_two_ports(pmi.pmi);
      case 2: return make_two_layer_two_ports(pmi.pmi);
      default: ocudu_assertion_failure("Unsupported PMI codebook configuration");
    }
    return {};
  }

  precoding_weight_matrix operator()(const pmi_typeI_single_panel& pmi) const
  {
    return make_type1_sp_mode1(pmi, nof_layers);
  }

  precoding_weight_matrix operator()(const pmi_typeII& pmi) const { return make_type2(pmi, nof_layers); }
};

precoding_weight_matrix ocudu::make_precoding(const precoding_matrix_indicator& pmi, unsigned nof_layers)
{
  return std::visit(precoding_calculator{nof_layers}, pmi);
}
\end{lstlisting}

Hier wird nichts gerechnet, nur verteilt.

\subsection{\file{lib/ran/precoding/type1\_sp/precoding\_codebooks.cpp}, die Umsetzung}

Die Implementierung von TS\,38.214 Abschnitt 5.2.2.2.1 für ein Single Panel:
\begin{itemize}
  \item \code{make\_layer\_type1\_sp\_mode1()}: ein Steuervektor $[v_{l,m};\,\varphi v_{l,m}]$ für einen Layer.
  \item \code{get\_k1\_and\_k2\_2\_layers()}, \code{get\_k1\_and\_k2\_3\_4\_layers()}: Tabellen 5.2.2.2.1-3 und -4.
  \item \code{make\_type1\_sp\_mode1()}: die vollständige Matrix Ports $\times$ Layer für 1 bis 8 Layer.
  \item \code{get\_beams\_from\_pmi()}, \code{calculate\_mimo\_matrix()}: die zerlegte Variante für Beam-Raster.
\end{itemize}

Der zugehörige Header \file{type1\_sp/precoding\_codebooks.h} liegt bewusst in \file{lib/} statt in
\file{include/}: Diese Funktionen sind intern, von außen soll nur über \code{make\_precoding()} zugegriffen werden.

\subsection{Warum die Aufteilung}

Der Rest des Systems soll nicht wissen, welcher Codebook-Typ gilt. Er ruft \code{make\_precoding(pmi, nof\_layers)}
und bekommt eine Matrix. Jeder Codebook-Typ hat sein eigenes Unterverzeichnis mit eigener Datei desselben Namens.
Zum Lesen geht man von \file{precoding\_codebooks.cpp:151} in den Visitor und von dort in
\file{type1\_sp/precoding\_codebooks.cpp}. Die dritte Datei, \file{type1\_sp/precoding\_codebook\_helpers.cpp},
liefert die Wertebereiche und Bitbreiten der PMI-Felder nach TS\,38.212 Tabelle 6.3.1.1.2-1, die der FAPI-Adaptor
für seine Schleifen und der CSI-Report-Decoder brauchen.

\section{Die Tests: Parametrisierung und Matrix-Ausgabe}

\subsection{Wer über die Topologien iteriert}

\file{tests/unittests/ran/beamforming/beam\_codebook\_generator\_test.cpp} enthält keine Schleife über die
Topologien; der Test wird sieben Mal als eigenständiger Testfall ausgeführt. Der Iterator ist Google Test.

\begin{lstlisting}[language=C++,title={tests/unittests/ran/beamforming/beam\_codebook\_generator\_test.cpp}]
class BeamCodebookGeneratorFixture : public ::testing::TestWithParam<antenna_topology>
{};

TEST_P(BeamCodebookGeneratorFixture, Generate)
{
  antenna_topology topology = GetParam();
  // ...
}

static constexpr auto antenna_topologies = to_array<antenna_topology>({antenna_topology::one_port,
                                                                       antenna_topology::two_port,
                                                                       antenna_topology::four_ports,
                                                                       antenna_topology::eight_ports,
                                                                       antenna_topology::single_panel_two_one,
                                                                       antenna_topology::single_panel_two_two,
                                                                       antenna_topology::single_panel_four_one});

INSTANTIATE_TEST_SUITE_P(BeamCodebookGenerator, BeamCodebookGeneratorFixture, ::testing::ValuesIn(antenna_topologies));
\end{lstlisting}

\begin{enumerate}
  \item Das Array ist nur die Datenquelle.
  \item \code{INSTANTIATE\_TEST\_SUITE\_P} ist der Auslöser. \code{ValuesIn()} verpackt das Array in einen
    \code{ParamGenerator<antenna\_topology>} (intern \code{ValuesInIteratorRangeGenerator}). Das Makro legt ein
    statisches Objekt an, das den Generator bei Programmstart in der gtest-Registry hinterlegt.
  \item Die Schleife steckt in gtest, \file{/usr/include/gtest/internal/gtest-param-util.h},
    \code{ParameterizedTestSuiteInfo::RegisterTests()}, ausgeführt bei \code{InitGoogleTest()}:
\begin{lstlisting}[language=C++]
for (typename ParamGenerator<ParamType>::iterator param_it = generator.begin();
     param_it != generator.end(); ++param_it, ++i) {
  // ...
  test_name_stream << test_info->test_base_name << "/" << param_name;
  MakeAndRegisterTestInfo(test_suite_name.c_str(), test_name_stream.GetString().c_str(), ...);
}
\end{lstlisting}
    Pro Wert entsteht ein Testfall \code{BeamCodebookGenerator/BeamCodebookGeneratorFixture.Generate/0} bis
    \code{/6}, jeder mit einer Factory, die den Parameterwert gespeichert hat.
  \item Zur Laufzeit erzeugt die Factory pro Testfall ein frisches Fixture-Objekt und setzt den Parameter.
    \code{GetParam()} liest ihn nur aus.
\end{enumerate}

\begin{lstlisting}[language=bash]
./cmake-build-debug/tests/unittests/ran/beam_codebook_generator_test --gtest_list_tests
./cmake-build-debug/tests/unittests/ran/beam_codebook_generator_test --gtest_filter='*/4'   # single_panel_two_one
\end{lstlisting}

\subsection{Beam-Gewichte als Matrix ausgeben}

Für \code{cf\_t} gibt es in \file{ocudu/adt/format.h} einen Formatter; \code{\{:+.2f\}} ergibt z.\,B.
\code{+0.50-0.50j}. Für \code{beam\_weights\_codebook} existiert kein Formatter, anders als für
\code{precoding\_weight\_matrix} (\file{include/ocudu/ran/precoding/precoding\_weight\_matrix\_formatters.h}).
Direkt im Test, eine Zeile pro Beam:

\begin{lstlisting}[language=C++]
fmt::print("Topologie {}: {} Beams x {} Antennen\n", to_string(topology), beam_weights.get_nof_beams(), nof_total_antennas);
for (unsigned i_beam = 0; i_beam != beam_weights.get_nof_beams(); ++i_beam) {
  fmt::print("  beam {:2d}: [", i_beam);
  for (unsigned i_antenna = 0; i_antenna != nof_total_antennas; ++i_antenna) {
    fmt::print("{}{:+.2f}", (i_antenna != 0) ? "  " : "",
               beam_weights.get_coefficient(static_cast<beam_identifier>(i_beam), i_antenna));
  }
  fmt::print("]\n");
}
\end{lstlisting}

Die erwartete Zeile im selben Layout, an der Stelle nach der Berechnung von \code{expected\_coefficients}:

\begin{lstlisting}[language=C++]
fmt::print("  expected beam {:2d} (panel={} pol={} l={} m={}): [",
           to_underlying(beam_id), i_panel, i_pol, i_beam_dim1, i_beam_dim2);
for (unsigned i_antenna = 0; i_antenna != nof_total_antennas; ++i_antenna) {
  fmt::print("{}{:+.2f}", (i_antenna != 0) ? "  " : "", expected_coefficients[i_antenna]);
}
fmt::print("]\n");
\end{lstlisting}

Die ersten Zeilen bis zur Antennenzahl sind die Einheitsmatrix (direkte Port-Beams), danach die DFT-Beams nach
Panel, Polarisation, $l$ und $m$. Die Zeilennummer ist die \code{beam\_identifier}.


\end{document}
